\section{Bounded sequence ingestion and integrity limits}
\label{sec:io-loading-contract}

Alignment and search operate on compact sequences, whereas public archives
usually distribute FASTA or FASTQ text. A FASTA record contains a header and a
sequence body. FASTQ adds a separator and a quality string whose length must
match the sequence. Records may span many input buffers, and one record may be
larger than a PostgreSQL value. Ingestion must therefore combine record
framing, compact encoding, metadata extraction, relational typing, and
transactional storage.

BioSeqDB separates these concerns into a reusable input layer and a
PostgreSQL-specific sink. The input layer reads bytes, selects the record
grammar, parses incrementally, and emits headers with compact sequence chunks.
It does not access PostgreSQL relations, tuples, or memory contexts. The sink
selects the representation declared by the destination sequence column,
converts header fields to SQL values, writes sequence fragments, and creates
the parent row that refers to them. It evaluates defaults and the executor row
constraints available to this bulk path. This boundary keeps file processing
independent of database internals while retaining transaction visibility and
rollback.

The resulting contract has four main properties. Multiline FASTA and FASTQ
records are framed across arbitrary refill boundaries. Sequence symbols are
packed directly into the destination representation. Packed sequence delivery
is bounded even for very long FASTA records. A failed load does not commit a
successful prefix as if the file had been loaded completely. FASTQ quality
storage has a narrower guarantee, stated in
Section~\ref{sec:io-integrity-limits}.

\section{Layered input pipeline}
\label{sec:io-reader-architecture}

The lowest layer follows a small byte-source contract. A source fills a
caller-owned span and reports the number of bytes read, end of input, or the
first error. The same parser can consequently consume a mapped plain file, a
sequential file descriptor, a zlib stream, or bytes supplied by another
provider. The parser is specialized for its source, so refill does not require
a virtual dispatch for every buffer.

Record format and compression are selected independently. The final
non-compression suffix selects the FASTA or FASTQ grammar. Common FASTA and
FASTQ suffixes are accepted case-insensitively, with an optional gzip suffix.
The opened file is then inspected for the gzip magic bytes defined by RFC
1952~\cite{rfc1952_gzip}. A compressed file can therefore be decompressed even
when its name omits a gzip suffix. A plain file with such a suffix is not sent
through a decompressor. By contrast, the parser does not guess FASTA or FASTQ
from the first record. An unrecognized or misleading format suffix produces a
grammar or format error instead of a heuristic reinterpretation.

For plain input, a mapped view is preferred when it is available and suitable
for the file size. Otherwise, sequential reads refill a reusable slab. The SQL
loader uses zlib for gzip input. Figure~\ref{fig:io-load-pipeline} shows where the
portable input work ends and database work begins.

\begin{figure}[htbp]
\centering
\begin{tikzpicture}[
    font=\small,
    box/.style={draw=black!65, rounded corners=2pt, align=center,
                minimum height=0.9cm, text width=2.55cm, fill=white},
    pg/.style={box, fill=blue!7},
    arr/.style={-{Latex[length=2mm]}, thick, draw=black!65},
    node distance=0.55cm and 0.55cm
]
  \node[box] (input) {FASTA or FASTQ file};
  \node[box, right=of input] (select) {select grammar and compression};
  \node[box, right=of select] (source) {mapped, sequential, or gzip source};
  \node[box, right=of source] (parser) {incremental record parser};

  \node[box, below=1.0cm of source] (pack) {alphabet-specific direct packing};
  \node[box, below=of pack] (chunks) {bounded packed chunks};
  \node[pg, below=of chunks] (fragments) {sequence fragments and locator};

  \node[box, below=1.0cm of parser] (header) {header dialect and field routing};
  \node[pg, below=of header] (typed) {SQL conversion, defaults, constraints};
  \node[pg, below=of typed] (parent) {parent row and locator publication};

  \draw[arr] (input) -- (select);
  \draw[arr] (select) -- (source);
  \draw[arr] (source) -- (parser);
  \draw[arr] (parser) -- (pack);
  \draw[arr] (pack) -- (chunks);
  \draw[arr] (chunks) -- (fragments);
  \draw[arr] (parser) -- (header);
  \draw[arr] (header) -- (typed);
  \draw[arr] (typed) -- (parent);
  \draw[arr] (fragments) -- (parent);
\end{tikzpicture}
\caption{Sequence and header data follow separate paths after record framing.
Packed sequence chunks are bounded. PostgreSQL operations are confined to the
shaded boxes.}
\label{fig:io-load-pipeline}
\end{figure}

\section{Incremental parsing and direct compact encoding}
\label{sec:io-parsing}

The parser is a finite-state machine over spans supplied by the selected
source. A streamed source refills a recyclable input slab. When a record
crosses a refill boundary, only its unread suffix is retained. A mapped source
can expose input bytes directly. The general batch reader may enlarge its slab
for an unusually large record, up to a configured hard limit. The SQL
streaming path instead drains complete packed units while the body is being
parsed. It does not first construct an ASCII sequence and then encode it, nor
does it require the complete packed body to be contiguous in memory.

For FASTA, the states distinguish the header marker, header text, sequence
body, and record end. Both LF and CRLF line endings are accepted, and the body
may contain arbitrary line wrapping. On the common DNA4 path, one vectorized
loop locates record boundaries and converts bases into two-bit storage words.
Other registered alphabets use the same state machine with their corresponding
packing operation. At record completion, the encoder flushes any partial word
and records its padding. The represented base count can then exclude unused
positions in the final storage unit.

FASTQ needs stricter framing because a quality line may begin with
\texttt{@}, \texttt{+}, or \texttt{>}. Searching for a marker at the start of
each physical line would therefore split valid records. The parser instead
performs the following steps:

\begin{enumerate}
    \item read one \texttt{@}-prefixed header
    \item collect one or more sequence lines until a \texttt{+}-prefixed
    separator
    \item skip the rest of the separator line
    \item collect quality bytes until their newline-free length equals the raw
    sequence length
\end{enumerate}

The final comparison uses the raw number of sequence bytes, not the number of
symbols retained by the alphabet encoder. This matters because the current
policy can discard an unknown symbol, while its corresponding quality byte is
still present. A truncated quality field is a terminal input error. Excess
quality bytes are rejected after a terminated line, but extra bytes in an
unterminated final line are currently accepted. Multiline sequence and quality
fields otherwise remain synchronized.

The destination sequence column controls representation. Before reading the
first record, the loader requires exactly one parameterized sequence column
and decodes its alphabet and storage-unit width. It then selects one compiled
combination of alphabet, unit width, and FASTA or FASTQ grammar. Available
layouts come from a compiled registry and can include DNA4, DNA5, and several
storage-unit widths. Unsupported combinations fail during schema inspection
rather than after a partial load. Each emitted chunk already uses the selected
compact layout, so the database sink appends its bytes without a second
conversion.

\section{Bounded delivery and backpressure}
\label{sec:io-bounds}

The SQL path emits a begin-record event, a sequence of packed chunks, and an
end-record event. Only complete storage units are emitted during a record. The
last chunk carries any partial unit and the adjusted base count. A chunk view
remains valid for its callback, allowing the database sink to copy it
directly into fragment storage. This protocol bounds the packed body by the
input slab and output window instead of the total sequence length. A
multi-gigabase FASTA sequence can thus be stored without creating one
multi-gigabyte database value.

PostgreSQL calls remain on the backend thread. The default path invokes the
sink directly, which gives deterministic order and avoids cross-thread
ownership. An optional overlap mode places owned begin, chunk, and end events
in a fixed-capacity first-in, first-out queue. A producer performs only
PostgreSQL-independent parsing. The backend thread consumes events and writes
relations. A full queue blocks the producer, which supplies explicit
backpressure. Closing the queue wakes both sides after an error, and the first
consumer failure remains the reported cause. The overlap mode is disabled by
default.

\section{Header-to-column mapping}
\label{sec:io-header-mapping}

Archive headers often combine an identifier with structured metadata.
UniProt, RefSeq, GenBank, NCBI, GISAID, and IMGT records use different
positional fields, delimiters, and tagged values. BioSeqDB can preserve the
entire header as text or extract recognized fields into SQL columns. The
parser supports an explicit dialect, a plain mode, and automatic detection.

Automatic detection is bounded and occurs once per input file. The loader
reads a limited decompressed prefix and samples a limited number of headers.
FASTQ sampling follows logical sequence and quality lengths, so an
\texttt{@} byte in quality data is not treated as a new header. Recognized
dialects vote across the sample with deterministic tie handling. If the sample
is unstructured and no metadata mapping requires further parsing, the loader
uses plain-header behavior. With metadata columns or user mappings, it retains
automatic mode and applies generic per-record parsing.

Mapping is also resolved once. Explicit user mappings take priority. Remaining
known fields are matched to normalized destination names, after which the
loader chooses an identifier column and routes any generic metadata. Name
normalization ignores ASCII case and common word separators. The resulting
routing table maps positional or tagged header slices to destination columns.
Per-record parsing then scans delimiters and passes raw byte ranges to the
database boundary without knowing PostgreSQL types.

At that boundary, each destination's input conversion and typmod are cached.
Unconstrained text can be constructed directly. Other scalar, range, or
multirange values use PostgreSQL's normal type input rules, including session
settings that affect dates and timestamps. Defaults are prepared once but
evaluated for every row, so volatile defaults retain their SQL meaning. The
completed tuple is checked with the executor's row-constraint step before it is
written. This direct bulk path does not execute ordinary \texttt{INSERT}
triggers or the secondary-index paths used for foreign-key, unique, and
exclusion enforcement. Per-row memory is reset only after referenced values
have been copied into database-owned storage.

The public operation can be summarized as
\texttt{bioseq.load(target\_table, file\_path, header\_format, columns)}. The
optional header format defaults to automatic detection, and the optional
mapping array can override inferred routes. The return value is the number of
loaded records. Ambiguous mappings, unsupported sequence layouts, generated
target columns, multiple sequence columns, and an unmapped required column
without a default are rejected before streaming begins.

\section{SQL publication and error behavior}
\label{sec:io-sql}

Sequence data is stored outside the small parent tuple. For each record, the
sink begins a sequence source, appends bounded packed chunks, and finishes the
source to obtain a pending locator. It then parses the header, evaluates
defaults, applies constraints, and writes the parent tuple containing that
locator. Publication occurs only after the parent row has copied the tuple.
This order avoids making an out-of-line sequence visible before a relational
row refers to it. Publication does not commit the row. Transaction visibility
still governs access to both objects.

The write is not equivalent to an ordinary \texttt{INSERT}. The loader writes
through its bulk storage path after evaluating defaults and executor row
constraints. It does not run normal insert triggers, maintain secondary
indexes, or invoke the usual foreign-key, unique, and exclusion paths. A target
relation must therefore be chosen with these limits in mind.

\begin{figure}[htbp]
\centering
\begin{minipage}{0.90\textwidth}
\small
\begin{verbatim}
inspect target and select format, alphabet, and unit width
detect the header dialect and resolve column routes

for each record
    begin an out-of-line sequence source
    append each bounded packed chunk
    finish the source and obtain a pending locator
    convert mapped header fields to SQL values
    evaluate unset defaults and enforce executor row constraints
    write the parent tuple with the locator
    publish the sequence source

flush parent and fragment writes
make the new rows visible to subsequent work in the transaction
optionally build configured sequence indexes
return the record count
\end{verbatim}
\end{minipage}
\caption{Simplified loading procedure. Representation and routing decisions
are made once, while row conversion and constraints are applied per record.}
\label{fig:io-load-pseudocode}
\end{figure}

Input and parser failures are carried to the PostgreSQL boundary with an owned
diagnostic. The boundary assigns an appropriate SQLSTATE and raises an SQL
error. PostgreSQL conversion, default, and constraint failures keep their
native messages. Error context adds the record number, the destination column
when known, and a bounded prefix of the header. Limiting the header excerpt
prevents a malformed record from creating an unbounded diagnostic.

Cleanup must preserve the earliest terminal error. Reader shutdown, queue
closure, and relation cleanup cannot replace it with a later secondary
failure. The target relation remains protected by a row-exclusive lock while
the loader owns its executor, tuple slot, and write states. On success, both
parent and fragment write states are flushed before newly loaded rows are used
for optional index construction. These resources are also released on the
error path. On failure, PostgreSQL aborts the enclosing transaction or
subtransaction. Rows and sequence fragments written before the fault are not
committed as a successful prefix. Within a larger explicit transaction, the
caller can roll back to a savepoint established before the load.

This design does not introduce a second transaction protocol in the input
layer. Locks, write-ahead logging, resource ownership, command visibility, and
rollback remain PostgreSQL responsibilities. The reusable parser reports
records and errors, while the database sink determines whether their effects
become visible.

\section{Integrity limits}
\label{sec:io-integrity-limits}

The bounded claim applies to packed FASTA sequence bodies. The FASTQ adapter
currently retains quality bytes until a record is complete, so its quality
scratch space grows with record length. Sequence packing is still incremental,
but total FASTQ memory is not independent of the largest record.

Alphabet checking is also permissive. Unknown FASTA and FASTQ symbols are
currently dropped by the encoder. A strict-alphabet option exists in the
reader interface but is not applied by the parser. The selected alphabet
therefore determines representation, not complete validation of every input
byte. Applications that require rejection of unknown symbols must validate
them before relying on the current loader.

The SQL gzip path uses zlib's stream validation. A decompression error after
output has begun is terminal because restarting with
another decoder would duplicate the already delivered prefix.

Finally, the filename suffix remains authoritative for choosing FASTA or
FASTQ. Magic-byte inspection identifies compression only. This separation
keeps parser dispatch deterministic, but a mislabeled biological format is
reported as an error rather than detected automatically.

\section{Conclusion}

BioSeqDB ingestion joins an incremental biological file parser to PostgreSQL
without materializing a complete large sequence. The parser frames records and
packs symbols into bounded native chunks. Header routing supplies typed column
values, and the database sink coordinates fragment storage with parent-row
publication and transaction rollback. The stated FASTQ and alphabet-validation
limits define where this guarantee does not yet apply.
